\vspace*{\fill}

\section{Introducción}

El presente proyecto desarrolla una plataforma web de punto de venta para \textit{Wallpa Wawqi}, orientada a mejorar la administraci\'on de productos, el registro de empleados y la presentaci\'on de informaci\'on del negocio de manera centralizada. La soluci\'on fue concebida como un sistema completo, compuesto por un backend en Java, un frontend desarrollado con Astro y una base de datos PostgreSQL, con el objetivo de cubrir tanto la l\'ogica de negocio como la experiencia de uso.

La necesidad principal del proyecto fue reemplazar procesos manuales y dispersos por una aplicaci\'on accesible desde la web, capaz de responder con rapidez a las operaciones habituales del negocio. Para ello, se defini\'o una arquitectura separada en capas, donde el backend expone servicios HTTP para autenticaci\'on y gesti\'on de productos, mientras que el frontend presenta una interfaz clara para el usuario final y consume dichos servicios de forma ordenada.

En el componente de servidor se emple\'o Java junto con clases organizadas por responsabilidad, incluyendo modelos de dominio, objetos de acceso a datos y servicios auxiliares para conexi\'on a la base de datos y almacenamiento de im\'agenes en la nube. En paralelo, el frontend fue dise\'nado para visualizar la informaci\'on principal del sistema, permitir el flujo de interacci\'on del usuario y mantener una estructura compatible con el despliegue en Vercel.

Adicionalmente, la persistencia de datos se resolvi\'o con PostgreSQL, alojada en Render.com, lo que permiti\'o mantener la informaci\'on disponible de forma estable para el backend desplegado en Railway. Esta combinaci\'on de tecnolog\'ias hizo posible construir una soluci\'on modular, funcional y f\'acil de mantener, alineada con los objetivos acad\'emicos del curso y con las necesidades reales del proyecto.

Finalmente, el documento presenta la problem\'atica identificada, la soluci\'on implementada, las conclusiones derivadas del desarrollo y las recomendaciones para una versi\'on futura. Las muestras, capturas de pantalla y demostraciones del funcionamiento del sistema se realizar\'an durante la exposici\'on, por lo que en esta versi\'on se han dejado espacios reservados para incorporar dichas evidencias.

\vspace*{\fill}
